% summary.tex — 5-page project summary: LLM-extracted carbon-reduction measures × Trucost emissions panel
\documentclass[11pt]{article}

\usepackage[margin=1in]{geometry}
\usepackage{booktabs}
\usepackage{amsmath}
\usepackage{amssymb}
\usepackage{graphicx}
\usepackage{array}
\usepackage{caption}
\usepackage[hidelinks]{hyperref}
\usepackage{microtype}

\captionsetup{font=small,labelfont=bf}
\setlength{\parskip}{0.25em}

\newcommand{\code}[1]{\texttt{\small #1}}

\title{\textbf{What Actually Reduces Corporate Carbon?}\\[0.3em]
\large A Two-Tier LLM Taxonomy of Sustainability-Report Measures\\
Linked to the Trucost Emissions Panel}
\author{Zheng Gong}
\date{\today}

\begin{document}
\maketitle
\vspace{-1.5em}

%======================================================================
\section{Research question and contribution}
%======================================================================

Firms publish long sustainability reports describing what they are doing about
climate change. Almost none of that text has been converted into
analysis-ready data. The literature instead studies either (i) \emph{disclosed
numbers} --- LLM pipelines that extract Scope~1/2/3 tonnages from PDFs, (ii)
\emph{disclosure quality} --- ClimateBERT/ChatReport-style scoring of TCFD
alignment and greenwashing, or (iii) \emph{panel regressions on CDP data},
which are bounded by CDP's voluntary respondents and CDP's own pre-defined
answer options. Nobody has taken the full text of the report, coded the
\emph{specific abatement actions} it describes into a standards-anchored
taxonomy, and asked whether those actions show up in realised emissions.

That is this project. The research question is deliberately blunt:

\begin{quote}
\textbf{Which carbon-reduction measures actually reduce carbon?} Conditional
on a firm adopting measure $m$, does the emission category that measure $m$
targets fall relative to the categories it does not target --- and which firm
characteristics predict both the \emph{choice} of measure and the
\emph{effectiveness} of that measure once chosen?
\end{quote}

Three things make the design work. First, the taxonomy is anchored in the GHG
Protocol scope architecture, so it is defensible to referees and, critically,
\emph{it maps one-to-one onto the outcome data}. Second, the outcome data
(S\&P Global Trucost, via WRDS) reports emissions at the same Scope~3 category
granularity that the GHG Protocol defines, which means a measure targeting
employee commuting has a matched outcome variable distinct from a measure
targeting use-phase emissions. Third, that matching delivers a
within-firm-year identification strategy: the untargeted categories of the
same firm in the same year act as the control group, absorbing everything
firm-specific --- growth, sector shocks, disclosure propensity, greenwashing
taste --- in a firm$\times$year fixed effect.

The deliverables are (a) a firm-year panel of $\approx$30 binary adoption
indicators for 1{,}892 US firms over 2015--2025, built from 11{,}756 report
PDFs; (b) estimates of the emission response to each measure, overall and by
targeted emission category; and (c) a map of which firm characteristics
operate through \emph{measure composition} versus \emph{measure efficacy}.

%======================================================================
\section{Measurement: the two-tier taxonomy}
%======================================================================

A single-level taxonomy forces a bad trade-off. Coarse categories throw away
the interesting variation; fine categories become unreliable when the LLM has
to read vague corporate prose. The taxonomy is therefore two-tiered, so that
extraction failure degrades gracefully rather than producing a missing
observation.

\textbf{Tier 1} is five high-level buckets tracking the accounting boundary of
the GHG Protocol: \textbf{S1} (direct emissions, Corporate Standard 2004),
\textbf{S2} (purchased energy, Scope~2 Guidance 2015), \textbf{S3U} (upstream
value chain, Categories 1--8), \textbf{S3D} (downstream value chain,
Categories 9--15), and \textbf{CDR} (removals and offsets, drawn from IPCC AR6
WGIII and the SBTi Corporate Net-Zero Standard). If a passage says only
``we reduced our operational footprint,'' Tier~1 still records S1 even though
no specific measure is identifiable.

\textbf{Tier 2} decomposes each bucket into 30 named measures. For Scope~3 the
structure is pre-given: the GHG Protocol fixes 15 categories, so the measure
list is not a judgement call. For Scopes~1 and 2 the measures are drawn from
IPCC AR6 WGIII Chapters 6, 9 and 11, the EU Taxonomy Regulation 2020/852, and
the Scope~2 Guidance. Two measures were added after a pilot revealed
systematic leakage (\code{renewable\_electricity\_general}, for renewable
power procured with no named instrument, and \code{packaging}); one Scope-3
category was split into supplier engagement versus material substitution
because the two are economically distinct and frequently co-occur.

\textbf{Governance} is treated as a parallel regressor set, not a Tier-1
bucket, because it is not a scope boundary: SBTi target validation, an
internal carbon price, executive compensation linked to emissions, and
third-party assurance. Keeping them separate lets us test whether governance
mechanisms predict realised abatement \emph{over and above} operational
measures --- a clean corporate-finance question in its own right.

\begin{table}[t]
\centering
\small
\caption{Taxonomy structure. 5 Tier-1 buckets, 30 Tier-2 measures, 4 governance flags.}
\begin{tabular}{@{}llp{9.2cm}@{}}
\toprule
\textbf{Tier 1} & \textbf{$n$} & \textbf{Tier-2 measures} \\
\midrule
S1  & 6 & energy efficiency; fuel switching / electrification; on-site renewables; F-gas substitution; methane \& fugitive capture; process emissions \\
S2  & 4 & PPA; REC/GoO; unspecified renewable electricity; low-carbon heat \& steam \\
S3U & 10 & c1 supplier engagement; c1 material substitution; c2 capital goods; c3 fuel \& energy; c4 upstream transport; c5 waste; c6 business travel; c7 commuting; c8 upstream leased; packaging \\
S3D & 7 & c9 downstream transport; c10 processing; c11 use phase; c12 end-of-life; c13 downstream leased; c14 franchises; c15 investments \\
CDR & 3 & nature-based removals; engineered CDR (DACCS/BECCS/biochar); voluntary offsets \\
\midrule
\multicolumn{3}{@{}l@{}}{\textbf{Governance (parallel):} SBTi target; internal carbon price; emission-linked executive pay; third-party assurance} \\
\bottomrule
\end{tabular}
\end{table}

The taxonomy lives in exactly one place in the code base
(\code{src/carbontax/taxonomy.py}); the LLM system prompt and the JSON
response schema are both \emph{derived} from that object rather than written
by hand, because prompt--schema drift is the dominant failure mode in this
kind of pipeline.

%======================================================================
\section{Extraction pipeline}
%======================================================================

The pipeline runs in four configuration-driven stages with no command-line
arguments; every knob is explicit in \code{config/run.yaml}.

\textbf{(1) Acquisition.} Sustainability-report PDFs are pulled from S\&P
Capital IQ for the Trucost company universe, giving a company$\leftrightarrow$filing
mapping of 13{,}273 filings. Of 13{,}177 filings requested for processing,
13{,}154 had a readable PDF (14 never downloaded, 9 corrupt).

\textbf{(2) Chunking.} Each PDF is parsed to text blocks with PyMuPDF,
repeated headers/footers are dropped by frequency, pages below 200 tokens are
discarded, and the remainder is split into 950-token windows with 50-token
overlap and \code{[PAGE~$n$]} markers preserved for traceability. A
case-insensitive keyword prefilter (\emph{carbon, emission, renewable, solar,
wind turbine, wind power, energy efficiency}) retains only climate-relevant
windows, and each filing is capped at 50 chunks by seeded random sampling.
Result: \textbf{177{,}041 chunks} from \textbf{11{,}756 filings} across
\textbf{1{,}892 companies} (median 10 chunks per filing, 35 per company).

\textbf{(3) Extraction.} One request per chunk to the OpenAI Batch API
(\code{gpt-5.4-nano}, temperature 0, strict JSON schema, prompt version
v2.2), sharded into 26 batches under the 20M enqueued-token cap. The schema is
deliberately flat --- 42 boolean properties, well under the 100-property
structured-output limit --- because a per-measure
\code{\{adopted, quote, page\}} object would have exceeded it. The prompt
imposes an \emph{action-not-mention} rule (``a measure is true only when the
chunk describes an action that reduces this source, not a mention that the
source exists'') and an \emph{aspiration} rule (``committed to\ldots'' and
``target to\ldots'' are false for measures, true for governance flags, since a
governance mechanism \emph{is} the commitment). Total input: 252.9M tokens,
$\approx$\$95 at batch-discounted rates --- i.e.\ \$0.05 per company, which is
what makes a wide-universe design feasible at all.

\textbf{(4) Aggregation.} Chunk-level flags are rolled up to the filing and
then the firm-year by taking the maximum: a measure is adopted if any chunk of
any report the firm filed that year evidences it. This yields
\textbf{8{,}090 firm-years} of adoption indicators.

%======================================================================
\section{Outcome data: Trucost Environmental (WRDS)}
%======================================================================

The dependent variables come from the WRDS \code{trucost\_environ} table
(21.1M rows, period end dates 2002--2026). The extract used here is the US
universe: \textbf{2{,}040 firms, 26{,}846 firm-years, FY2002--2026}, with
Scope~1, Scope~2 location-based, Scope~3 upstream total and Trucost revenue
populated for 99.9\% of rows. Merging the extraction panel to Trucost on
\code{companyid} $\times$ fiscal year retains \textbf{7{,}944} firm-years
contemporaneously and \textbf{7{,}864} at a one-year lag --- essentially no
attrition, because the report universe was drawn from the Trucost universe by
construction.

The decisive feature is the category-level Scope~3 detail, which lines up with
the Tier-2 taxonomy almost measure-for-measure (Table~\ref{tab:map}).

\begin{table}[t]
\centering
\small
\caption{Matched action$\rightarrow$outcome pairs. This mapping is what makes the
within-firm-year design possible.}
\label{tab:map}
\begin{tabular}{@{}llll@{}}
\toprule
\textbf{Tier-2 measure} & \textbf{Trucost outcome} & \textbf{Variable} & \textbf{Coverage} \\
\midrule
S1 measures (6)            & Absolute GHG Scope 1                & \code{di\_319413} / \code{di\_376883} & 99.9\% \\
S2 measures (4)            & Scope 2 location-based              & \code{di\_319414} / \code{di\_376884} & 99.9\% \\
\quad PPA, REC/GoO         & Scope 2 \emph{market-based}         & \code{di\_367750} / \code{di\_376886} & 9.5\% \\
c1 supplier / material     & Purchased goods \& services         & \code{di\_378459} & as-reported \\
c2 capital goods           & Capital goods                       & \code{di\_378680} & as-reported \\
c3 fuel \& energy          & Fuel \& energy services             & \code{di\_378681} & as-reported \\
c4 upstream transport, c6 travel & Upstream T\&D + business travel (pooled) & \code{di\_378682} & as-reported \\
c5 waste                   & Waste generated                     & \code{di\_378683} & as-reported \\
c7 commuting               & Employee commuting                  & \code{di\_378684} & as-reported \\
c8 upstream leased         & Upstream leased assets              & \code{di\_378685} & as-reported \\
c9 downstream transport    & Downstream T\&D                     & \code{di\_368652} & as-reported \\
c10 processing             & Processing of sold products         & \code{di\_368653} & as-reported \\
c11 use phase              & Use of sold products                & \code{di\_368654} & as-reported \\
c12 end-of-life            & EOL treatment of sold products      & \code{di\_368655} & as-reported \\
c13 downstream leased      & Downstream leased assets            & \code{di\_368704} & as-reported \\
c14 franchises             & Franchises                          & \code{di\_368705} & as-reported \\
c15 investments            & Investments (financed emissions)    & \code{di\_368706} & as-reported \\
CDR measures (3)           & \emph{no inventory counterpart}     & ---               & --- \\
\midrule
Scale / deflator           & Trucost total revenue               & \code{di\_319522} & 99.9\% \\
Disclosure quality         & Weighted disclosure: GHG, Environmental & \code{di\_319416}, \code{di\_329704} & 99.9\% \\
Per-scope disclosure score & GHG Scope 1/2/3 scores              & \code{di\_387688}--\code{di\_387726} & --- \\
\bottomrule
\end{tabular}
\end{table}

Three constraints follow directly and shape the empirical design. (i) The
theoretically correct outcome for PPAs and RECs is \emph{market-based} Scope~2,
which is populated for only 9.5\% of firm-years --- so instrument-specific S2
measures are estimable only on a much smaller sub-panel, while
location-based S2 is the right outcome for on-site generation and efficiency.
(ii) Trucost pools upstream transport with business travel, so c4 and c6
cannot be separated on the outcome side even though the taxonomy separates
them on the action side. (iii) CDR has no inventory counterpart by
construction --- removals sit outside the GHG inventory --- which turns the
CDR flags into a natural placebo and a natural greenwashing test rather than a
treatment with a matched outcome.

%======================================================================
\section{Descriptive evidence}
%======================================================================

Table~\ref{tab:adopt} reports firm-year adoption rates. Two patterns matter
for the regressions. First, the distribution is extremely skewed: energy
efficiency appears in 73.6\% of firm-years while five Scope-3 categories
(c2, c8, c10, c13, c14) appear in under 1\%, so those measures carry no
estimable variation and should be pruned or pooled. The mean firm-year
reports 4.2 distinct measures (median 4, p90 = 9). Second, adoption is
strongly trending --- SBTi targets rise from 5.6\% (2016) to 27.5\% (2023),
RECs from 11.9\% to 21.7\%, PPAs from 6.0\% to 13.7\%, engineered CDR from
1.0\% to 5.1\% --- while energy efficiency is flat at $\approx$74\%. Any
specification must therefore absorb year effects, and preferably
sector$\times$year effects.

\begin{table}[t]
\centering
\small
\caption{Firm-year adoption rates (\% of 8{,}090 firm-years), selected measures.}
\label{tab:adopt}
\begin{tabular}{@{}lr lr lr@{}}
\toprule
\multicolumn{2}{c}{\textbf{Tier 1}} & \multicolumn{2}{c}{\textbf{Tier 2 (top)}} & \multicolumn{2}{c}{\textbf{Tier 2 (thin) \& governance}} \\
\cmidrule(r){1-2}\cmidrule(lr){3-4}\cmidrule(l){5-6}
S1  & 78.9 & Energy efficiency        & 73.6 & Business travel (c6)    & 6.4 \\
S2  & 64.4 & Supplier engagement (c1) & 39.1 & Internal carbon price   & 5.9 \\
S3U & 58.2 & On-site renewables       & 38.0 & Downstream transport (c9) & 4.7 \\
CDR & 26.7 & Waste / circularity (c5) & 30.3 & Engineered CDR          & 3.1 \\
S3D & 24.9 & Renewable elec.\ (generic) & 28.0 & Investments (c15)     & 2.6 \\
    &      & Packaging                & 25.3 & Fuel \& energy (c3)     & 1.1 \\
    &      & Material substitution (c1) & 21.8 & Capital goods (c2)    & 0.5 \\
    &      & Third-party assurance    & 36.9 & Downstream leased (c13) & 0.4 \\
    &      & SBTi target              & 19.1 & Franchises (c14)        & 0.4 \\
    &      & Exec.\ pay linked        & 17.1 & Upstream leased (c8)    & 0.4 \\
    &      & Fuel switching           & 18.2 & Processing (c10)        & 0.3 \\
\bottomrule
\end{tabular}
\end{table}

A third descriptive result is a finding in its own right: the
\emph{unpinpointed share} --- chunks where Tier~1 fires but no Tier-2 measure
does --- is 23.7\% for S1, 15.0\% for S3U, 16.5\% for CDR, but 41.1\% for S3D
and 54.4\% for S2. Over half of all Scope-2 talk names no instrument at all.
That is a direct, quantified measure of how much corporate climate disclosure
is scope-level assertion rather than identifiable action, and it is also why
the generic \code{renewable\_electricity\_general} category had to exist.

%======================================================================
\section{Empirical design (current stage)}
%======================================================================

\paragraph{Specification 1 --- firm-level baseline.} For firm $i$ in sector
$s$ and year $t$,
\begin{equation}
\Delta \ln E_{it} \;=\; \sum_{m=1}^{M}\beta_m D_{i,m,t-1}
\;+\; \gamma' X_{it} \;+\; \alpha_i \;+\; \delta_{st} \;+\; \varepsilon_{it},
\end{equation}
where $E_{it}$ is Scope~1$+$2 absolute emissions (log intensity, deflated by
Trucost revenue, as the alternative outcome), $D_{i,m,t-1}$ are the lagged
adoption indicators, $\alpha_i$ are firm fixed effects, and $\delta_{st}$ are
sector$\times$year effects. $X_{it}$ must include \emph{disclosure volume} ---
the number of retained chunks and the Trucost weighted-disclosure score ---
because a firm that simply writes a longer report mechanically trips more
flags. Without that control, $\beta_m$ confounds doing with talking. Standard
errors are two-way clustered by firm and year.

\paragraph{Specification 2 --- matched category panel (the identifying design).}
Stack the data at the firm$\times$category$\times$year level using
Table~\ref{tab:map}, where $k$ indexes an emission category and $D^{\text{match}}_{ikt}$
equals one if the firm adopted a measure targeting category $k$:
\begin{equation}
\Delta \ln E_{ikt} \;=\; \beta\, D^{\text{match}}_{i,k,t-1}
\;+\; \alpha_{it} \;+\; \gamma_{kt} \;+\; \varepsilon_{ikt}.
\end{equation}
The firm$\times$year fixed effect $\alpha_{it}$ absorbs \emph{every}
firm-level confounder in that year: output growth, M\&A, sector shocks,
disclosure propensity, general greenwashing taste, even the firm's overall
climate ambition. Identification comes purely from the contrast between the
categories a firm targeted and the categories it did not, within the same
report. Interacting $D^{\text{match}}$ with category dummies recovers a
measure-specific $\beta_k$ under the same fixed-effect structure, which is
precisely the ``which action works for which emission type'' question.
The design also generates its own falsification test: the coefficient on
$D_{i,k',t-1}$ for a \emph{non}-matched category $k' \neq k$ should be zero,
and the CDR measures --- which have no inventory counterpart --- should move
nothing.

\paragraph{Specification 3 --- dynamics and staggered adoption.} Because
adoption is staggered and trending, a static two-way fixed-effects estimator
is contaminated by already-treated comparisons. Estimate an event study around
each firm's \emph{first} adoption of measure $m$, $\tau \in [-3,+5]$, using
not-yet-treated firms as controls and a heterogeneity-robust estimator
(Callaway--Sant'Anna or Sun--Abraham). Pre-trends are the key diagnostic here:
firms plausibly announce measures \emph{after} emissions have already begun to
fall, which is exactly the anticipation problem the event study is meant to
expose.

\paragraph{Sample discipline.} Trucost \emph{models} emissions from
revenue $\times$ sector intensity when a firm does not disclose. On modelled
observations, emissions are a mechanical function of revenue and cannot
respond to a firm's abatement measures --- so including them attenuates every
$\beta$ toward zero. The primary sample must therefore be restricted to
as-reported observations using the \code{*\_As Reported} variables and the
per-scope GHG disclosure scores (\code{di\_387688}--\code{di\_387726}), with
the modelled observations used as an explicit placebo: an effect that also
appears in the modelled subsample is a spurious revenue artefact.

%======================================================================
\section{Next stage: firm characteristics, composition, and efficacy}
%======================================================================

The natural next question is not just \emph{which} measures work but
\emph{for whom}, and through what channel. Let $Z_{it}$ collect firm
characteristics: size (log assets), Tobin's $q$, leverage, ROA, cash holdings,
capex and R\&D intensity, asset tangibility, firm age, institutional ownership
and its concentration, index membership, analyst coverage, financial
constraints (SA/WW indices), product-market competition, foreign sales share,
baseline emission intensity relative to sector, and policy exposure (state RPS
stringency, regulated utility status, carbon-price jurisdiction share).

\paragraph{(a) Moderated effects.} Interact characteristics with measures
inside Specification 2:
\begin{equation}
\Delta \ln E_{ikt} = \big(\beta + \lambda' Z_{i,t-1}\big) D^{\text{match}}_{i,k,t-1}
+ \alpha_{it} + \gamma_{kt} + \varepsilon_{ikt},
\end{equation}
where the level effect of $Z$ is absorbed by $\alpha_{it}$ --- a real
advantage, since it means $\lambda$ is identified off differential
\emph{effectiveness} rather than off any correlation between firm type and
emission growth.

\paragraph{(b) Composition versus efficacy.} A characteristic can raise
realised abatement in two distinct ways: by shifting firms toward measures
that work (composition), or by making the same measure more effective
(efficacy). A Oaxaca--Blinder-style decomposition of the gap between firm
types $z_1$ and $z_0$ separates them:
\begin{equation}
\underbrace{\textstyle\sum_m \bar{\beta}_m \big[p_m(z_1)-p_m(z_0)\big]}_{\text{composition: chooses better actions}}
\;+\;
\underbrace{\textstyle\sum_m \bar{p}_m \big[\beta_m(z_1)-\beta_m(z_0)\big]}_{\text{efficacy: executes actions better}},
\end{equation}
with $p_m(z)$ estimated from a first-stage adoption model (a multivariate
probit or LPM of $D_m$ on $Z$, which is independently interesting: it is the
selection equation describing \emph{who does what}). The finance interpretation
is sharp --- if slack-rich, low-constraint firms decarbonise mainly through
composition, capital availability buys \emph{options}; if through efficacy, it
buys \emph{execution}.

\paragraph{(c) Governance as a moderator, not a treatment.} The four
governance flags are best read here as candidate moderators. Does an internal
carbon price make efficiency investment more effective, or merely more
frequent? Does emission-linked executive pay raise $\beta$, or does it raise
adoption of cheap, visible measures (offsets, RECs) without moving the
inventory? Third-party assurance is doubly useful: it is both a plausible
moderator and a proxy for data quality, so it can enter as a sample split.

\paragraph{(d) Flexible heterogeneity.} With $\approx$30 treatments and
$\approx$20 characteristics, the interaction space is too large to explore by
hand without inviting specification search. A causal-forest / generic-ML
approach (Chernozhukov--Demirer--Duflo--Fern\'andez-Val: GATES and CLAN)
delivers pre-registered, sample-split heterogeneity estimates --- discipline
that will matter to referees given how many interactions are available.

%======================================================================
\section{Threats to validity}
%======================================================================

\textbf{Measurement error in an LLM-generated regressor} is the first-order
concern. Classical error in a binary regressor attenuates $\beta$; correlated
error (a model systematically over-firing on firms with glossier reports) does
worse. Mitigation: a hand-coded validation sample with per-measure precision
and recall, repeated extraction under different models and seeds to bound
instability, and where feasible the recent econometrics for regressors
generated by machine learning. \textbf{Report--fiscal-year alignment}: a
report filed in year $t$ typically describes fiscal year $t-1$, so the lag
structure must be verified against \code{periodenddate} rather than assumed.
\textbf{The 50-chunk cap} censors long reports, so adoption is measured with
more noise for the most disclosure-heavy firms --- estimate a robustness
sample without the cap. \textbf{Selection into disclosure} is the deepest
issue: firms choose to publish a report, choose what to say, and choose
whether to report emissions to Trucost. Specification 2 defends against the
firm-level component of this but not against the possibility that firms
selectively disclose measures for the categories where they already expect
improvement --- which is precisely what the event-study pre-trends are for.

%======================================================================
\section{Future directions}
%======================================================================

\textbf{1. From binary to dosage.} The current flags record whether a measure
was mentioned as an action, not how large it was. A second extraction pass can
pull quantities --- MW of contracted PPA capacity, dollars of abatement capex,
percentage targets, share of renewable electricity --- turning $D_m$ into a
continuous treatment and greatly sharpening the effect estimates.

\textbf{2. The say--do gap as an asset-pricing object.} Firms whose measure
adoption predicts \emph{no} subsequent emission response are, operationally,
the definition of cheap talk. A firm-level ``say--do gap'' index --- predicted
minus realised abatement --- can be validated against known greenwashing
events and then priced: do high-gap firms earn lower returns, face higher
credit spreads, or suffer larger drawdowns on climate-news days? This turns a
measurement project into a finance paper.

\textbf{3. Sharper identification.} Candidate shocks include state RPS
tightenings, the SEC climate-disclosure rule and California SB 253/261, the EU
CBAM for exporters, and utility-level fuel-price shocks. Each provides
plausibly exogenous variation in the \emph{cost} of a specific measure family,
supporting instrumented or difference-in-differences estimates of measure
effectiveness rather than the conditional-correlation reading of
Specification~1.

\textbf{4. Supply-chain spillovers.} Supplier engagement (c1) is the most
common Scope-3 measure. Its effect should appear not only in the focal firm's
Scope~3 but in \emph{suppliers'} Scope~1. Linking the panel to a supply-chain
network (FactSet Revere or Compustat customer segments) tests whether
value-chain pressure propagates, and whether it merely relocates emissions
outside the reporting boundary.

\textbf{5. Offsets versus abatement.} Because CDR has no inventory
counterpart, firms leaning on offsets should show adoption without inventory
response. Comparing offset-heavy against measure-heavy decarbonisers speaks
directly to the SBTi's insistence that removals cannot substitute for
reduction.

\textbf{6. Open-vocabulary discovery.} The unpinpointed residual (up to 54\%
of Scope-2 passages) is where the fixed taxonomy fails. Embedding-clustering
that residual text would surface measure types the standards do not enumerate,
allowing the taxonomy to be extended empirically rather than only deductively.

\textbf{7. Scope extension.} The pipeline is universe-agnostic and costs
\$0.05 per company. Extending beyond the US --- where the EU's CSRD now
mandates comparable disclosure --- would allow a regulatory-regime comparison:
does mandated disclosure change what firms do, or only what they write?

\vspace{1em}
\begin{center}
\rule{0.4\textwidth}{0.4pt}
\end{center}
\small
\noindent\textbf{Data and code.} Extraction pipeline: \code{carbontax}
(Python, editable install; four stages driven by \code{config/run.yaml}).
Reports: S\&P Capital IQ ESG filings. Emissions: S\&P Global Trucost
Environmental via WRDS (\code{trucost\_environ}). Financials: Compustat
North America. Standards: GHG Protocol Corporate Standard (2004), Scope~2
Guidance (2015), Corporate Value Chain (Scope~3) Standard (2011); IPCC AR6
WGIII (2022); EU Taxonomy Regulation 2020/852; SBTi Corporate Net-Zero
Standard v1.3 (2025).

\end{document}
